\begin{proofbox}
\paragraph{方法一：直接求导检验误差（优先读，无须高阶展开）}
这条路线先验证给定分式，再给出可用于舍入的误差界。

\textbf{1. 指数：比较 $x$ 与 $\ln R_E(x)$。}
令 $P=x^2+6x+12$，$Q=x^2-6x+12$。
由 $P=(x+3)^2+3>0$、$Q=(x-3)^2+3>0$，可取对数。设
\[
H(x)=x-\ln\frac PQ,\qquad H(0)=0.
\]
只求一次导数并通分：
\[
H'(x)=1-\frac{P'}P+\frac{Q'}Q
=\frac{x^4}{PQ}=\frac{x^4}{x^4-12x^2+144}\ge0.
\]
因此 $H$ 在 $x>0$ 时为正、$x<0$ 时为负，得到
\[
x>0:\ R_E(x)<e^x;\qquad x<0:\ R_E(x)>e^x.
\]
这里的严格性来自 $H'(x)>0$（$x\ne0$），并非由图像猜出。

当 $0\le x\le1$，令 $u=x^2\in[0,1]$。由于 $u^2-12u+144$ 在此区间递减，$PQ\ge1-12+144=133$。
所以 $0\le H'(x)\le x^4/133$。比较 $H$ 与 $x^5/665$ 的导数，并使用它们在零点的值都为零，得
\[
0\le H(x)\le\frac{x^5}{665}.
\]
对指数函数用中值定理，且 $\ln R_E\le x\le1$：
\[
0\le e^x-R_E(x)=e^x-e^{\ln R_E(x)}
\le e\,H(x)<\frac{3x^5}{665}\quad(x>0).
\]
这就证明了结论中的误差界。这里使用常数估计 $e<3$。

\medskip
\textbf{2. 对数：直接对“原函数减分式”求导。}
令 $D=x^2+6x+6$，$K(x)=\ln(1+x)-R_L(x)$。当 $x>-1$ 时 $D>1$，且
\[
K(0)=0,\qquad K'(x)=\frac{x^4}{(1+x)D^2}\ge0.
\]
故 $x>0$ 时 $K(x)>0$，$-1<x<0$ 时 $K(x)<0$。
在 $0\le x\le1$ 上，$1+x\ge1$、$D\ge6$，于是
\[
0\le K'(x)\le\frac{x^4}{36}.
\]
比较 $K$ 与 $x^5/180$ 的导数和零点值，便得
\[
0\le K(x)\le\frac{x^5}{180}.
\]
这一证明只用求导、通分和单调性；也说明为什么不能把近似式不加证明地用于放缩。

\newpage
\paragraph{方法二：配系数，解释公式为什么长这样}
这一角度需知道零点附近的泰勒展开；它用于\textbf{构造}分式，并逐一验证五个公式。
将分母常数项归一化为 $1$，交叉相乘后只比较低次项，无须直接对分式反复求高阶导数。

\textbf{1. 指数。}设 $R=(1+ax+bx^2)/(1+cx+dx^2)$，要求
\[
(1+cx+dx^2)e^x-(1+ax+bx^2)=O(x^5).
\]
由 $e^x=1+x+x^2/2+x^3/6+x^4/24+O(x^5)$，三、四次项给出
\[
\frac16+\frac c2+d=0,\qquad
\frac1{24}+\frac c6+\frac d2=0.
\]
解得 $c=-1/2,d=1/12$；一次、二次项给出 $a=1/2,b=1/12$，即 $R=R_E$。
分母在零点非零，交叉相乘的余项阶数不会因除以分母而改变。

\textbf{2. 对数。}设 $R=(ax+bx^2)/(1+cx+dx^2)$，用
\[
\ln(1+x)=x-\frac{x^2}2+\frac{x^3}3-\frac{x^4}4+O(x^5).
\]
令三、四次项消失：
\[
\frac13-\frac c2+d=0,\qquad
-\frac14+\frac c3-\frac d2=0.
\]
解得 $c=1,d=1/6$，再得 $a=1,b=1/2$，便是 $R_L$。

\textbf{3. 正弦与正切。}先利用奇对称，设 $R=x/(1+bx^2)$。
因为
\[
\frac{x}{1+bx^2}=x-bx^3+O(x^5),
\]
与 $\sin x=x-x^3/6+O(x^5)$ 比较得 $b=1/6$，即 $R_S$；
与 $\tan x=x+x^3/3+O(x^5)$ 比较得 $b=-1/3$，即 $R_T$。

\textbf{4. 余弦。}先利用偶对称，设 $R=(1+ax^2)/(1+bx^2)$。展开为
\[
R=1+(a-b)x^2+(b^2-ab)x^4+O(x^6).
\]
与 $\cos x=1-x^2/2+x^4/24+O(x^6)$ 比较，得
\[
a-b=-\frac12,\qquad b(b-a)=\frac1{24}.
\]
所以 $b=1/12,a=-5/12$，得到 $R_C$。

\medskip
\textbf{首个误差项（均指 $x\to0$）：}
\[
\begin{aligned}
e^x-R_E(x)&=\frac{x^5}{720}+O(x^6),&
\ln(1+x)-R_L(x)&=\frac{x^5}{180}+O(x^6),\\
\sin x-R_S(x)&=-\frac{7x^5}{360}+O(x^7),&
\cos x-R_C(x)&=\frac{x^6}{480}+O(x^8),\\
\tan x-R_T(x)&=\frac{x^5}{45}+O(x^7).
\end{aligned}
\]
这些系数由继续比较下一非零项得到。它们解释误差图的趋势；若要保证具体小数位，仍须像方法一那样控制余项。

\newpage
\paragraph{数形结合：奇偶性来自单位圆}
\begin{center}
\begin{tikzpicture}[font=\small]
\draw[black!30,rounded corners=2pt] (0,0) rectangle (14.4,8.2);
\node[anchor=west,font=\small\bfseries] at (0.5,7.72) {单位圆：横坐标不变，纵坐标反号};

\draw[-Stealth] (1.8,4)--(8.5,4) node[right] {$u$};
\draw[-Stealth] (5,1)--(5,6.85) node[above] {$v$};
\draw[black!60] (5,4) circle (2.5);
\draw (7.5,3.94)--(7.5,4.06);
\node[anchor=west] at (7.65,4.25) {$1$};
\coordinate (P) at ({5+2.5*cos(30)},{4+2.5*sin(30)});
\coordinate (Q) at ({5+2.5*cos(30)},{4-2.5*sin(30)});
\draw[blue,very thick] (5,4)--(P);
\draw[orange!85!black,thick,dashed] (5,4)--(Q);
\draw[densely dotted] (P)--(Q);
\fill[blue] (P) circle(2pt);
\fill[orange!85!black] (Q) circle(2pt);
\draw[-Stealth] (6.05,4) arc[start angle=0,end angle=30,radius=1.05];
\draw[-Stealth] (6.05,4) arc[start angle=0,end angle=-30,radius=1.05];
\node at (6.35,4.27) {$\alpha$};
\node at (6.35,3.73) {$-\alpha$};
\node[anchor=south west,blue] at (7.5,5.45) {$P=(\cos\alpha,\sin\alpha)$};
\node[anchor=north west,orange!85!black] at (7.5,2.55) {$Q=(\cos\alpha,-\sin\alpha)$};
\node[anchor=west,align=left] at (9,4) {$\cos(-\alpha)=\cos\alpha$\\[6pt]$\sin(-\alpha)=-\sin\alpha$};
\node[anchor=north east] at (4.85,3.9) {$O$};
\node[anchor=west] at (0.6,0.55) {图取 $\alpha=\pi/6=30^\circ$；对应短弧长为 $\pi/6$（半径为 $1$）。};

\end{tikzpicture}
\end{center}

角 $\alpha$ 与 $-\alpha$ 对应的两点关于横轴对称，故横坐标相同、纵坐标相反。这给出余弦偶、正弦奇，解释了方法二为何选只含偶次幂的余弦分式、奇函数型的正弦分式。
图取 $\alpha=\pi/6$；一般对称关系对任意实数角度都成立。角的数值用于展开和代入时必须转为弧度。

\paragraph{方法三：对数的面积推导（选读）}
此方法需接受“$\ln(1+x)$ 是 $y=1/(1+t)$ 在 $[0,x]$ 下方的面积”这一微积分事实，以下设 $x>0$。
\begin{center}
\begin{tikzpicture}[font=\small]
\draw[black!30,rounded corners=2pt] (0,0) rectangle (14.4,8.2);
\node[anchor=west,font=\small\bfseries] at (0.5,7.72) {以 $x=3$ 为例：对数分式来自两次高度采样};
\node[anchor=west] at (0.5,6.95) {原曲线下面积};
\begin{axis}[width=5.2cm,height=5.3cm,at={(1.85cm,1.25cm)},anchor=south west,axis lines=left,xmin=0,xmax=3.2,ymin=0,ymax=1.15,xlabel={$t$},ylabel={高度},xtick={0,1.5,3},ytick={0,0.5,1}]
\addplot[draw=none,fill=blue!12,domain=0:3,samples=100]{1/(1+x)} \closedcycle;
\addplot[blue,very thick,domain=0:3,samples=100]{1/(1+x)};
\end{axis}
\node[anchor=west] at (7.5,6.95) {两矩形近似面积};
\begin{axis}[width=5.2cm,height=5.3cm,at={(8.25cm,1.25cm)},anchor=south west,axis lines=left,xmin=0,xmax=3.2,ymin=0,ymax=1.15,xlabel={$t$},ylabel={高度},xtick={0,1.5,3},ytick={0,0.5,1}]
\path[fill=orange!20] (axis cs:0,0) rectangle (axis cs:1.5,0.6120046189);
\path[fill=orange!20] (axis cs:1.5,0) rectangle (axis cs:3,0.2970862902);
\draw[orange!85!black,thick] (axis cs:0,0)--(axis cs:0,0.6120046189)--(axis cs:1.5,0.6120046189)--(axis cs:1.5,0.2970862902)--(axis cs:3,0.2970862902)--(axis cs:3,0);
\addplot[only marks,orange!85!black] coordinates {(0.6339745962,0.6120046189) (2.3660254038,0.2970862902)};
\node[anchor=south west] at (axis cs:0.65,0.73) {$t_-$};
\node[anchor=south] at (axis cs:2.3,0.42) {$t_+$};
\addplot[blue,very thick,domain=0:3,samples=100]{1/(1+x)};
\end{axis}

\end{tikzpicture}
\end{center}

以两个宽均为 $x/2$ 的矩形近似面积，在
\[
t_-=\frac x2-\frac{x}{2\sqrt3},\qquad
 t_+=\frac x2+\frac{x}{2\sqrt3}
\]
处取曲线高度。两矩形总面积为
\[
A_2=\frac x2\left(\frac1{1+t_-}+\frac1{1+t_+}\right)
=\frac{x(1+x/2)}{(1+x/2)^2-x^2/12}
=\frac{6x+3x^2}{6+6x+x^2}=R_L(x).
\]
\textbf{采样点为何不是各半段的普通中点？}设两采样点为 $x/2\pm b$。
对常数和一次函数，两矩形法已给出准确面积；要使二次函数也准确，须
\[
\frac x2\bigl[(x/2-b)^2+(x/2+b)^2\bigr]
=\int_0^x t^2\,\mathrm dt=\frac{x^3}3,
\]
故 $b^2=x^2/12$。三次函数也准确，因为
\[
\frac x2\bigl[(x/2-b)^3+(x/2+b)^3\bigr]=\frac{x^4}4.
\]
因此，这不是随手挑两个点，而是一种对三次及以下多项式准确的面积近似。
还可用只需多项式除法的精确恒等式说明误差阶数：
\[
\frac1{1+t}=1-t+t^2-t^3+\frac{t^4}{1+t}.
\]
在 $0\le t\le x$ 上，最后一项介于 $0$ 与 $x^4$ 之间，所以它在真实面积及两矩形面积中的贡献均不超过 $x^5$。
前三次多项式的面积完全匹配，故总面积误差不超过 $2x^5$，特别是 $O(x^5)$。这个界较粗，但无须高阶导数。
这给出同一个对数分式的几何来源；它的严格上下关系与数值界已由方法一证明。
\end{proofbox}
